% Change these commands to alter every rendered cross-reference and inline format.
\newcommand{\CharacterReferenceStyle}[1]{\textcolor{wine}{\textbf{#1}}}
\newcommand{\SpaceReferenceStyle}[1]{\textit{#1}}
\newcommand{\ClueReferenceStyle}[1]{\textbf{#1}}
\newcommand{\MysteryReferenceStyle}[1]{\textcolor{wine}{\textbf{#1}}}
\newcommand{\TimestampStyle}[1]{\texttt{#1}}
\newcommand{\OutsideEntityStyle}[1]{\textbf{#1}}
\newcommand{\MainObjectiveStyle}[1]{\textbf{#1}}
\newcommand{\PositiveRewardStyle}[1]{\textcolor{green}{#1}}
\newcommand{\NegativeRewardStyle}[1]{\textcolor{red}{#1}}

% Semantic heading styles.  Their defaults deliberately preserve LaTeX's
% standard heading appearance; customize each independently when needed.
\newcommand{\ObjectivesHeadingStyle}[1]{#1}
\newcommand{\LosingConditionsHeadingStyle}[1]{#1}
\newcommand{\SecretsHeadingStyle}[1]{#1}
\newcommand{\ActionsHeadingStyle}[1]{#1}
\newcommand{\BackgroundHeadingStyle}[1]{#1}
\newcommand{\DescriptionHeadingStyle}[1]{#1}

% Page background configuration.  Set this to an image path relative to the
% repository root (for example, content/assets/background.png).  Leave it
% empty to keep page backgrounds disabled.
\newcommand{\BackgroundImagePath}{content/assets/background/dark-grunge-background.jpg}
\newcommand{\DefaultBackgroundOpacity}{0.15}
